\newcommand{\tablaER}[4]{%
  \node[entidad] (#1) at #2 {\textcolor{azul}{\textbf{\texttt{#3}}}\\[2pt]
    \scalebox{0.72}{\begin{tabular}{@{}l@{\hspace{3pt}}l@{\hspace{3pt}}l@{}}#4\end{tabular}}};}
\begin{frame}[label=er-completo]{Ítem 4 · Diagrama entidad-relación completo de PostgreSQL}
\centering
{\scriptsize 10 tablas · 47 columnas · 8 claves foráneas · Vista de consulta derivada}\par\vspace{2pt}
\begingroup
\begin{tikzpicture}[
  entidad/.style={draw=azul!45,fill=azul!3,rounded corners=2pt,anchor=north west,
    inner sep=4pt,text width=3.12cm,align=left,font=\tiny},
  relacion/.style={draw=verde!80!black,line width=0.65pt,-{Latex[length=1.2mm]}},
  every label/.style={font=\tiny}]
\tablaER{tiempo}{(0,0)}{dim\_tiempo}{
anio & INTEGER & PK\\
periodo & VARCHAR(20) & *\\
es\_pandemia & BOOLEAN & *\\}
\tablaER{edad}{(0,-1.28)}{dim\_rango\_edad}{
rango\_edad & VARCHAR(20) & PK\\
grupo\_etario & VARCHAR(25) & *\\
orden & SMALLINT & *\\
edad\_min & SMALLINT & *\\
edad\_max & SMALLINT & *\\}
\tablaER{sexo}{(0,-2.98)}{dim\_sexo}{
sexo & VARCHAR(15) & PK\\
descripcion & VARCHAR(25) & *\\}
\tablaER{educativo}{(0,-4.10)}{dim\_nivel\_educativo}{
nivel\_educativo & VARCHAR(25) & PK\\
orden & SMALLINT & *\\}
\tablaER{hecho}{(3.62,0)}{hecho\_desaparicion}{
id & BIGSERIAL & PK\\
anio & INTEGER & FK*\\
mes & SMALLINT & *\\
ubigeo & CHAR(6) & FK\\
sexo & VARCHAR(15) & FK*\\
rango\_edad & VARCHAR(20) & FK*\\
estado\_civil & VARCHAR(25) & FK*\\
nacionalidad & VARCHAR(60) & FK*\\
nivel\_educativo & VARCHAR(25) & FK*\\
ocupacion & VARCHAR(80) & FK*\\
cantidad & SMALLINT & *\\}
\tablaER{civil}{(3.62,-3.42)}{dim\_estado\_civil}{
estado\_civil & VARCHAR(25) & PK\\}
\tablaER{ocupacion}{(3.62,-4.40)}{dim\_ocupacion}{
ocupacion & VARCHAR(80) & PK\\}
\tablaER{geografia}{(7.24,0)}{dim\_geografia}{
ubigeo & CHAR(6) & PK\\
departamento & VARCHAR(60) & *\\
provincia & VARCHAR(60) & *\\
distrito & VARCHAR(80) & *\\
region\_natural & VARCHAR(10) & *\\
cod\_departamento & CHAR(2) & *\\
cod\_provincia & CHAR(4) & *\\}
\tablaER{nacionalidad}{(7.24,-2.16)}{dim\_nacionalidad}{
nacionalidad & VARCHAR(60) & PK\\
es\_extranjera & BOOLEAN & *\\}
\node[entidad,fill=verde!5,draw=verde!50] (reglas) at (7.24,-3.45) {
\textbf{Integridad y relaciones}\\[3pt]
Relación \textbf{1:N}: dimensión a hechos.\\
Las FK usan la PK homónima.\\
\texttt{mes}: entre 1 y 12.\\
\texttt{cantidad}: mayor que 0.\\
\texttt{region\_natural}: COSTA, SIERRA o SELVA.\\
\texttt{ubigeo} del hecho admite NULL.};
\tablaER{staging}{(10.86,0)}{stg\_desaparicion}{
anio & INTEGER & \\
mes & INTEGER & \\
ubigeo\_hecho & VARCHAR(10) & \\
dpto\_hecho & VARCHAR(60) & \\
prov\_hecho & VARCHAR(60) & \\
dist\_hecho & VARCHAR(80) & \\
sexo & VARCHAR(15) & \\
rango\_edad & VARCHAR(20) & \\
est\_civil & VARCHAR(25) & \\
nacionalidad & VARCHAR(60) & \\
nivel\_educativo & VARCHAR(25) & \\
ocupacion & VARCHAR(80) & \\
cantidad & INTEGER & \\}
\node[entidad,draw=naranja!60,fill=naranja!4] (vista) at (10.86,-3.65) {
\textbf{\texttt{v\_desaparicion} · VISTA}\\[2pt]
\texttt{anio, periodo, mes}\\
\texttt{departamento, provincia}\\
\texttt{distrito, region\_natural}\\
\texttt{sexo, grupo\_etario}\\
\texttt{rango\_edad, nacionalidad}\\
\texttt{es\_extranjera, nivel\_educativo}\\
\texttt{estado\_civil, ocupacion}\\
\texttt{cantidad}\\
Derivada de hechos + tiempo,\\geografía, edad y nacionalidad.};
\draw[relacion] (tiempo.east) -- (hecho.west |- tiempo.east);
\draw[relacion] (edad.east) -- (hecho.west |- edad.east);
\draw[relacion] (sexo.east) -- ++(0.12,0) |- ($(hecho.north west)!0.55!(hecho.south west)$);
\draw[relacion] (educativo.east) -- ++(0.20,0) |- ($(hecho.north west)!0.82!(hecho.south west)$);
\draw[relacion] (geografia.west) -- (hecho.east |- geografia.west);
\draw[relacion] (nacionalidad.west) -- ++(-0.12,0) |- ($(hecho.north east)!0.72!(hecho.south east)$);
\draw[relacion] (civil.north) -- (hecho.south);
\draw[relacion] (ocupacion.west) -- ++(-0.15,0) |- ([yshift=0.08cm]hecho.south west);
\end{tikzpicture}
\endgroup
\par{\tiny
PK: clave primaria (única y no nula). FK: clave foránea. *: NOT NULL.
Staging recibe el CSV mediante ETL; no tiene PK/FK.
\enlace{etl}{Volver al proceso ETL}. Fuente: \texttt{sql/01\_ddl.sql}.}
\end{frame}
